% TOP_PROBLEM: {"id":30007662,"problem_number":"LOCAL-30007662","title":"Predicting institutional turning points","category_id":21,"category":{"id":21,"name":"economics","display_name":"Economics","description":"Open economic theory statements, published research agendas, and proposed scoped specifications; question provenance and historical priority are qualified per record.","slug":"economics","order_index":21,"created_at":"2026-10-08T00:00:00Z"},"status":"research_question","external_url":null,"legacy_ids":[],"published":false,"statement":"In the explicitly specified setting: Can researchers identify critical junctures in real time and explain which beliefs and coordination processes produce durable institutional change? The mathematical statement above fixes the target, assumptions and scope of this catalogue formulation.","economics":{"original_id":151,"field":"Political economy and institutions","kind":"identification","formalization_basis":"adapted_research_question","source_status":"research_agenda","priority_status":"earliest_found","priority_note":"Earliest directly inspected qualifying passage in this audit; earlier origins were not established. A review or research-agenda statement does not imply original priority.","earliest_verified_reference":{"authors":["Michael Callen","Jonathan L. Weigel","Noam Yuchtman"],"title":"Experiments about Institutions","year":2023,"url":"https://www.ifo.de/DocDL/cesifo1_wp10833.pdf","doi":"10.1146/annurev-economics-091823-031317","locator":"CESifo WP 10833, section 3.1, printed p. 7; section 4.2, printed p. 13","evidence_summary":"The authors call for systematic real-time beliefs measurement and improved identification of institutional critical junctures."},"current_reference":{"authors":["Michael Callen","Jonathan L. Weigel","Noam Yuchtman"],"title":"Experiments about Institutions","year":2023,"url":"https://www.ifo.de/DocDL/cesifo1_wp10833.pdf","doi":"10.1146/annurev-economics-091823-031317","locator":"CESifo WP 10833, section 3.1, printed p. 7; section 4.2, printed p. 13","evidence_summary":"The authors call for systematic real-time beliefs measurement and improved identification of institutional critical junctures."},"formalization_note":"The cited passage establishes a research direction, not this exact mathematical statement. The notation, population, policy menu, assumptions, and estimands are an original adaptation for this catalogue; no universal unsolved theorem or source priority is claimed."},"statusQualification":"Published question or research agenda verified at the cited locator. The mathematical specification is adapted for this catalogue and includes additional assumptions; source publication does not establish first-ever priority or certify this exact model remains unsolved. The cited passage establishes a research direction, not this exact mathematical statement. The notation, population, policy menu, assumptions, and estimands are an original adaptation for this catalogue; no universal unsolved theorem or source priority is claimed.","statusReviewedAt":"2026-10-08"}
% FIRST_POSTED: 2026-10-08
% ENTRY_KIND: statement_only
% !TeX program = lualatex
\documentclass[11pt]{article}
\usepackage{fontspec}
\usepackage[a4paper,margin=25mm]{geometry}
\usepackage{amsmath,amssymb,mathtools}
\usepackage{xurl}
\usepackage[hidelinks]{hyperref}
\newcommand{\sref}[2]{\hyperlink{source:#1:#2}{[S#2]}}
\setlength{\emergencystretch}{3em}
\begin{document}
% problemId:30007662
\section*{Predicting institutional turning points}\label{30007662}
\subsection{Definitions and mathematical statement}
Fix a panel of polities with institutional states \(R_{ct}\) measured by a declared rule at date \(t\). Citizen \(i\) in polity \(c\) reports a probability distribution \(b_{ict}(r)\) over possible institutions \(r\) at horizon \(T\). Define a prespecified uncertainty index \(H_{ct}\) using within-person entropy and between-person disagreement, and let \(X_{ct}\) contain information observed before outcomes. The prediction task is to estimate \(\Pr(R_{cT}\ne R_{ct}\mid H_{ct},X_{ct})\) and test calibration in episodes selected before institutional changes are known. A critical juncture may exist even when the status quo survives, so samples cannot contain only successful transitions. For a feasible binary information or coordination intervention \(z\in\{0,1\}\), define the causal estimand \(\tau_R=E[1\{R_{cT}(1)\in\mathcal I\}-1\{R_{cT}(0)\in\mathcal I\}]\), where \(\mathcal I\) is a declared set of inclusive institutional states. Distinguish predictive uncertainty from uncertainty causing collective action. Measure higher-order beliefs where relevant, address survey selection and repression, and predefine durability. The cited authors propose real-time belief measurement; this quantified forecasting and intervention problem adapts their agenda and does not claim a universal sufficient statistic for institutional change.

\par\textbf{Formulation and scope.} The cited passage establishes a research direction, not this exact mathematical statement. The notation, population, policy menu, assumptions, and estimands are an original adaptation for this catalogue; no universal unsolved theorem or source priority is claimed.

\par\textbf{Open-question provenance.} An explicit question or research agenda was verified at the source locator. This does not by itself certify that every later version remains unresolved \sref{30007662}{1}.

\par\textbf{Historical priority.} Earliest directly inspected qualifying passage in this audit; earlier origins were not established. A review or research-agenda statement does not imply original priority.

\subsection{Short English statement}
In the explicitly specified setting: Can researchers identify critical junctures in real time and explain which beliefs and coordination processes produce durable institutional change? The mathematical statement above fixes the target, assumptions and scope of this catalogue formulation.

\subsection{Sources}
\begin{itemize}\raggedright
\item[\textnormal{[S1]}]\hypertarget{source:30007662:1}{} Michael Callen, Jonathan L. Weigel, Noam Yuchtman. 2023. Experiments about Institutions. CESifo WP 10833, section 3.1, printed p. 7; section 4.2, printed p. 13.
\url{https://www.ifo.de/DocDL/cesifo1_wp10833.pdf}.
DOI: 10.1146/annurev-economics-091823-031317.
{\small\textit{Earliest verified question/agenda reference; historical priority qualified below. The authors call for systematic real-time beliefs measurement and improved identification of institutional critical junctures.}}
\end{itemize}
\end{document}
